% p115_1; source PDF page 115; manually verified editable table.
\begin{tikzpicture}[x=1cm,y=-1cm,inner sep=0pt,outer sep=0pt]
\path[use as bounding box] (0,0) rectangle (16.000,3.165);
\draw[line width=.4pt] (0.000,0.000) rectangle (4.600,0.855);
\node[align=center,text width=4.340cm,font={\linespread{1}\fontsize{10}{13.50}\selectfont\bfseries}] at (2.300,0.427) {种类};
\draw[line width=.4pt] (4.600,0.000) rectangle (9.300,0.855);
\node[align=center,text width=4.440cm,font={\linespread{1}\fontsize{10}{13.50}\selectfont\bfseries}] at (6.950,0.427) {认识对象};
\draw[line width=.4pt] (9.300,0.000) rectangle (16.000,0.855);
\node[align=center,text width=6.440cm,font={\linespread{1}\fontsize{10}{13.50}\selectfont\bfseries}] at (12.650,0.427) {法律后果};
\draw[line width=.4pt] (0.000,0.855) rectangle (4.600,2.160);
\node[align=left,text width=4.340cm,font={\linespread{1}\fontsize{10}{13.50}\selectfont}] at (2.300,1.508) {事实认识错误};
\draw[line width=.4pt] (4.600,0.855) rectangle (9.300,2.160);
\node[align=left,text width=4.440cm,font={\linespread{1}\fontsize{10}{13.50}\selectfont}] at (6.950,1.508) {构成要件事实（例如行为对象）};
\draw[line width=.4pt] (9.300,0.855) rectangle (16.000,2.160);
\node[align=left,text width=6.440cm,font={\linespread{1}\fontsize{10}{13.50}\selectfont}] at (12.650,1.508) {排除犯罪故意};
\draw[line width=.4pt] (0.000,2.160) rectangle (4.600,3.165);
\node[align=left,text width=4.340cm,font={\linespread{1}\fontsize{10}{13.50}\selectfont}] at (2.300,2.663) {法律认识错误};
\draw[line width=.4pt] (4.600,2.160) rectangle (9.300,3.165);
\node[align=left,text width=4.440cm,font={\linespread{1}\fontsize{10}{13.50}\selectfont}] at (6.950,2.663) {刑法的禁止性（违法性）};
\draw[line width=.4pt] (9.300,2.160) rectangle (16.000,3.165);
\node[align=left,text width=6.440cm,font={\linespread{1}\fontsize{10}{13.50}\selectfont}] at (12.650,2.663) {若无违法性认识可能性，则排除责任};
\end{tikzpicture}
